% GENERATED FILE. Edit canonical lessons, then run npm run generate:books.
% canonical-source-hash: fnv1a64:043ebb388ae69664
% canonical-lessons: TE-C257-vagdanam, TE-C257-rahasyam, TE-C257-sneham, TE-C257-godava, TE-C257-samskrti, TE-R257-first-pass-words-from-chapters-249-252, TE-R257-second-pass-words-from-chapters-253-257

\chapter{Promise, Secret, Friendship, Quarrel, Culture}
\label{ch:te-promise-secret-friendship-quarrel-culture}
\hlchaptermodality{257}

\begin{chapteropening}
\textbf{By the end of this chapter:} \emph{I can say a promise, a secret, friendship, a quarrel and culture.}

Complete the last lesson of chapter 257: I can say a promise, a secret, friendship, a quarrel and culture.
\end{chapteropening}

\section[\texorpdfstring{\te{వాగ్దానం} (vāgdānaṁ)}{వాగ్దానం (vāgdānaṁ)}]{\te{వాగ్దానం} (vāgdānaṁ) — a promise}

\label{lesson:TE-C257-vagdanam}



\begin{quote}
Before the new one: say the Telugu for a warehouse, then the Telugu for a court.
\end{quote}

\subsection*{You'll want to know: \te{వాగ్దానం}}
\textbf{\te{వాగ్దానం}} — \emph{vāgdānaṁ} — \textquotedblleft{}a promise\textquotedblright{}.

\begin{grammarlens}[title={What you've built}]
One more word for this chapter.
\end{grammarlens}

\subsection*{Your turn}
\begin{itemize}
\raggedright
  \item \emph{Say it:} \emph{vāgdānaṁ}
  \item \emph{Say it:} \emph{vāgdānaṁ}, once more
  \item \emph{Say it:} say \emph{vāgdānaṁ}
  \item \emph{Recall:} say the Telugu for a warehouse, then the Telugu for a court, then say \emph{vāgdānaṁ} again
\end{itemize}

\begin{tcolorbox}[breakable,colback=teal!4,colframe=teal!35!black,title={Before you move on}]
What does \te{వాగ్దానం} mean? (\textquotedblleft{}A promise\textquotedblright{}.) Say it once more.
\end{tcolorbox}
\section[\texorpdfstring{\te{రహస్యం} (rahasyaṁ)}{రహస్యం (rahasyaṁ)}]{\te{రహస్యం} (rahasyaṁ) — a secret}

\label{lesson:TE-C257-rahasyam}



\begin{quote}
Before the new one: say the Telugu for a court, then the Telugu for a promise.
\end{quote}

\subsection*{You'll want to know: \te{రహస్యం}}
\textbf{\te{రహస్యం}} — \emph{rahasyaṁ} — \textquotedblleft{}a secret\textquotedblright{}.

\begin{grammarlens}[title={What you've built}]
One more word for this chapter.
\end{grammarlens}

\subsection*{Your turn}
\begin{itemize}
\raggedright
  \item \emph{Say it:} \emph{rahasyaṁ}
  \item \emph{Say it:} \emph{rahasyaṁ}, once more
  \item \emph{Say it:} say \emph{rahasyaṁ}
  \item \emph{Recall:} say the Telugu for a court, then the Telugu for a promise, then say \emph{rahasyaṁ} again
\end{itemize}

\begin{tcolorbox}[breakable,colback=teal!4,colframe=teal!35!black,title={Before you move on}]
What does \te{రహస్యం} mean? (\textquotedblleft{}A secret\textquotedblright{}.) Say it once more.
\end{tcolorbox}
\section[\texorpdfstring{\te{స్నేహం} (snēhaṁ)}{స్నేహం (snēhaṁ)}]{\te{స్నేహం} (snēhaṁ) — friendship}

\label{lesson:TE-C257-sneham}



\begin{quote}
Before the new one: say the Telugu for a promise, then the Telugu for a secret.
\end{quote}

\subsection*{You'll want to know: \te{స్నేహం}}
\textbf{\te{స్నేహం}} — \emph{snēhaṁ} — \textquotedblleft{}friendship\textquotedblright{}.

\begin{grammarlens}[title={What you've built}]
One more word for this chapter.
\end{grammarlens}

\subsection*{Your turn}
\begin{itemize}
\raggedright
  \item \emph{Say it:} \emph{snēhaṁ}
  \item \emph{Say it:} \emph{snēhaṁ}, once more
  \item \emph{Say it:} say \emph{snēhaṁ}
  \item \emph{Recall:} say the Telugu for a promise, then the Telugu for a secret, then say \emph{snēhaṁ} again
\end{itemize}

\begin{tcolorbox}[breakable,colback=teal!4,colframe=teal!35!black,title={Before you move on}]
What does \te{స్నేహం} mean? (\textquotedblleft{}Friendship\textquotedblright{}.) Say it once more.
\end{tcolorbox}
\section[\texorpdfstring{\te{గొడవ} (goḍava)}{గొడవ (goḍava)}]{\te{గొడవ} (goḍava) — a quarrel, a fuss}

\label{lesson:TE-C257-godava}



\begin{quote}
Before the new one: say the Telugu for a secret, then the Telugu for friendship.
\end{quote}

\subsection*{You'll want to know: \te{గొడవ}}
\textbf{\te{గొడవ}} — \emph{goḍava} — \textquotedblleft{}a quarrel, a fuss\textquotedblright{}.

\begin{grammarlens}[title={What you've built}]
One more word for this chapter.
\end{grammarlens}

\subsection*{Your turn}
\begin{itemize}
\raggedright
  \item \emph{Say it:} \emph{goḍava}
  \item \emph{Say it:} \emph{goḍava}, once more
  \item \emph{Say it:} say \emph{goḍava}
  \item \emph{Recall:} say the Telugu for a secret, then the Telugu for friendship, then say \emph{goḍava} again
\end{itemize}

\begin{tcolorbox}[breakable,colback=teal!4,colframe=teal!35!black,title={Before you move on}]
What does \te{గొడవ} mean? (\textquotedblleft{}A quarrel, a fuss\textquotedblright{}.) Say it once more.
\end{tcolorbox}
\section[\texorpdfstring{\te{సంస్కృతి} (saṁskṛti)}{సంస్కృతి (saṁskṛti)}]{\te{సంస్కృతి} (saṁskṛti) — culture}

\label{lesson:TE-C257-samskrti}



\begin{quote}
Before the new one: say the Telugu for friendship, then the Telugu for a quarrel.
\end{quote}

\subsection*{You'll want to know: \te{సంస్కృతి}}
\textbf{\te{సంస్కృతి}} — \emph{saṁskṛti} — \textquotedblleft{}culture\textquotedblright{}.

\begin{grammarlens}[title={What you've built}]
One more word for this chapter.
\end{grammarlens}

\subsection*{Your turn}
\begin{itemize}
\raggedright
  \item \emph{Say it:} \emph{saṁskṛti}
  \item \emph{Say it:} \emph{saṁskṛti}, once more
  \item \emph{Say it:} say \emph{saṁskṛti}
  \item \emph{Recall:} say the Telugu for friendship, then the Telugu for a quarrel, then say \emph{saṁskṛti} again
\end{itemize}

\begin{tcolorbox}[breakable,colback=teal!4,colframe=teal!35!black,title={Before you move on}]
What does \te{సంస్కృతి} mean? (\textquotedblleft{}Culture\textquotedblright{}.) Say it once more.
\end{tcolorbox}
\section[(dialogue)]{Words from chapters 249-252 — a first pass}

\label{lesson:TE-R257-first-pass-words-from-chapters-249-252}



\begin{quote}
Say the words for a quarrel and culture.
\end{quote}

\subsection*{Your turn}
Say these aloud:

\begin{itemize}
\raggedright
  \item to endure, to crawl, to bend down, to stretch out, to stack
  \item the palm of the hand, the heel, the ankle, the thigh, the waist
  \item the mind, breathlessness, a muscle, an injection, a vaccination
  \item Ugadi, Diwali, Dasara, Vinayaka Chavithi, Bathukamma
\end{itemize}

\begin{tcolorbox}[breakable,colback=teal!4,colframe=teal!35!black,title={Before you move on}]
To endure? (\textbf{\te{భరించు}}, \emph{bhariñcu}.) To crawl? (\textbf{\te{పాకు}}, \emph{pāku}.) To bend down? (\textbf{\te{వంగు}}, \emph{vaṅgu}.) To stretch out? (\textbf{\te{చాచు}}, \emph{cācu}.) To stack? (\textbf{\te{పేర్చు}}, \emph{pērcu}.) The palm of the hand? (\textbf{\te{అరచెయ్యి}}, \emph{araceyyi}.) The heel? (\textbf{\te{మడమ}}, \emph{maḍama}.) The ankle? (\textbf{\te{చీలమండ}}, \emph{cīlamaṇḍa}.) The thigh? (\textbf{\te{తొడ}}, \emph{toḍa}.) The waist? (\textbf{\te{నడుము}}, \emph{naḍumu}.) The mind? (\textbf{\te{మనసు}}, \emph{manasu}.) Breathlessness? (\textbf{\te{ఆయాసం}}, \emph{āyāsaṁ}.) A muscle? (\textbf{\te{కండరం}}, \emph{kaṇḍaraṁ}.) An injection? (\textbf{\te{ఇంజెక్షను}}, \emph{iñjekṣanu}.) A vaccination? (\textbf{\te{టీకా}}, \emph{ṭīkā}.) Ugadi? (\textbf{\te{ఉగాది}}, \emph{ugādi}.) Diwali? (\textbf{\te{దీపావళి}}, \emph{dīpāvaḷi}.) Dasara? (\textbf{\te{దసరా}}, \emph{dasarā}.) Vinayaka Chavithi? (\textbf{\te{వినాయక} \te{చవితి}}, \emph{vināyaka caviti}.) Bathukamma? (\textbf{\te{బతుకమ్మ}}, \emph{batukamma}.)
\end{tcolorbox}
\section[(dialogue)]{Words from chapters 253-257 — a second pass}

\label{lesson:TE-R257-second-pass-words-from-chapters-253-257}



\begin{quote}
Say the words for the end and ever.
\end{quote}

\subsection*{Your turn}
Say these aloud:

\begin{itemize}
\raggedright
  \item the end, ever, at present, in the past, dawn
  \item a state, an area, a direction, east, west
  \item a museum, a gym, a stage, a city, the village meeting platform
  \item a lane, a tunnel, the outskirts, a warehouse, a court
  \item a promise, a secret, friendship, a quarrel, culture
\end{itemize}

\begin{tcolorbox}[breakable,colback=teal!4,colframe=teal!35!black,title={Before you move on}]
The end? (\textbf{\te{చివర}}, \emph{civara}.) Ever? (\textbf{\te{ఎప్పుడైనా}}, \emph{eppuḍainā}.) At present? (\textbf{\te{ప్రస్తుతం}}, \emph{prastutaṁ}.) In the past? (\textbf{\te{పూర్వం}}, \emph{pūrvaṁ}.) Dawn? (\textbf{\te{తెల్లవారుజాము}}, \emph{tellavārujāmu}.) A state? (\textbf{\te{రాష్ట్రం}}, \emph{rāṣṭraṁ}.) An area? (\textbf{\te{ప్రాంతం}}, \emph{prāntaṁ}.) A direction? (\textbf{\te{దిక్కు}}, \emph{dikku}.) East? (\textbf{\te{తూర్పు}}, \emph{tūrpu}.) West? (\textbf{\te{పడమర}}, \emph{paḍamara}.) A museum? (\textbf{\te{మ్యూజియం}}, \emph{myūziyaṁ}.) A gym? (\textbf{\te{వ్యాయామశాల}}, \emph{vyāyāmaśāla}.) A stage? (\textbf{\te{వేదిక}}, \emph{vēdika}.) A city? (\textbf{\te{నగరం}}, \emph{nagaraṁ}.) The village meeting platform? (\textbf{\te{రచ్చబండ}}, \emph{raccabaṇḍa}.) A lane? (\textbf{\te{సందు}}, \emph{sandu}.) A tunnel? (\textbf{\te{సొరంగం}}, \emph{soraṅgaṁ}.) The outskirts? (\textbf{\te{శివారు}}, \emph{śivāru}.) A warehouse? (\textbf{\te{గోదాము}}, \emph{gōdāmu}.) A court? (\textbf{\te{కోర్టు}}, \emph{kōrṭu}.) A promise? (\textbf{\te{వాగ్దానం}}, \emph{vāgdānaṁ}.) A secret? (\textbf{\te{రహస్యం}}, \emph{rahasyaṁ}.) Friendship? (\textbf{\te{స్నేహం}}, \emph{snēhaṁ}.) A quarrel? (\textbf{\te{గొడవ}}, \emph{goḍava}.) Culture? (\textbf{\te{సంస్కృతి}}, \emph{saṁskṛti}.)
\end{tcolorbox}
